\documentclass[addpoints,spanish, 12pt,a4paper]{exam}
\pointpoints{punto}{puntos}
\hpword{Puntos:}
\vpword{Puntos:}
\htword{Total}
\vtword{Total}
\hsword{Resultado:}
\hqword{Ejercicio:}
\vqword{Ejercicio:}

\printanswers

\usepackage[utf8]{inputenc}
\usepackage[spanish]{babel}
\usepackage{eurosym}
\usepackage{yhmath}
%\usepackage[spanish,es-lcroman, es-tabla, es-noshorthands]{babel}

\usepackage{verbatim}

\usepackage[margin=1in]{geometry}
\usepackage{amsmath,amssymb}
\usepackage{multicol}

\usepackage{graphicx}
\graphicspath{{../img/}} 

\newcommand{\class}{Matemáticas 4º Aplicadas}
\newcommand{\examdate}{\today}
\newcommand{\examnum}{Expresiones algebraicas}
\newcommand{\tipo}{A}


\newcommand{\timelimit}{50 minutos}

\renewcommand{\solutiontitle}{\noindent\textbf{Solución:}\enspace}

\pagestyle{head}
\firstpageheader{\includegraphics[width=0.2\columnwidth]{header_left}}{\textbf{Departamento de Matemáticas\linebreak \class}\linebreak \examnum}{\includegraphics[width=0.1\columnwidth]{header_right}}
\runningheader{\class}{\examnum}{Página \thepage\ of \numpages}
\runningheadrule

\pointsinrightmargin % Para poner las puntuaciones a la derecha. Se puede cambiar. Si se comenta, sale a la izquierda.
\extrawidth{-2.4cm} %Un poquito más de margen por si ponemos textos largos.
\marginpointname{ \emph{\points}}


\begin{document}

\noindent
\begin{tabular*}{\textwidth}{l @{\extracolsep{\fill}} r @{\extracolsep{6pt}} }
\textbf{Nombre:} \makebox[3.5in]{\hrulefill} & \textbf{Fecha:}\makebox[1in]{\hrulefill} \\
 & \\
\textbf{Tiempo: \timelimit} & Tipo: \tipo 
\end{tabular*}
\rule[2ex]{\textwidth}{2pt}
Esta prueba tiene \numquestions\ ejercicios. La puntuación máxima es de \numpoints. 
La nota final de la prueba será la parte proporcional de la puntuación obtenida sobre la puntuación máxima. Para la evaluación de pendientes de 3ºESO o 2ºPMAR se tendrán en cuenta los apartados 1, 2, 3 y 6.a y 6.c: 

\begin{center}


\addpoints
 %\gradetable[h][questions]
	\pointtable[h][questions]
\end{center}

\noindent
\rule[2ex]{\textwidth}{2pt}

\begin{questions}

\question Indica el coeficiente, la parte literal y el grado de los siguientes monomios.
 
\begin{parts}
\part [1] $A=6x^3$

\begin{solution}$ $\end{solution}
\part [1] $B=-3x$

\begin{solution}$ $\end{solution}
\part [1] $C=4x^3$

\begin{solution}$ $\end{solution}
\end{parts}

\addpoints

\question Dados los monomios $A=6x^3$, $B=-3x$, $C=4x^3$, calcula:
\begin{parts}
\part [1] $(C-A)\cdot B$

\begin{solution}$ 6x^4$\end{solution}
\part [1] $\dfrac{B\cdot C}{A}$

\begin{solution}$-2x$\end{solution}

\part [1] $\dfrac{A^2}{2B}$
\begin{solution}$-6x^5 $\end{solution}

\end{parts}
\addpoints


\question [2] Opera y simplifica la siguiente expresión: $$\left(2x^2-3x+1\right)\left(2x-1\right)-\left(4x^3-8x^2+1\right) $$
\begin{solution}
$5x-2$
\end{solution}

\question [2] Multiplica por 6 esta expresión y simplifica: $$\dfrac{2x^2-1}{2}-\dfrac{x-1}{3}-\dfrac{1-x}{6}$$
\begin{solution}
$ 6x^2-x-2$
\end{solution}

\question  Resuelve las siguientes ecuaciones:
\begin{parts}
\part [1] $3(2x+1)-\dfrac{x+1}{2}=\dfrac{1}{2}(x+2-\dfrac{x+1}{3}) $

\begin{solution}$\frac{-10}{31} $\end{solution}
\part [1] $\dfrac{2x+1}{5}-\dfrac{x+1}{3}=\dfrac{3x}{5}-2(\dfrac{x}{6}-1)$

\begin{solution} $\frac{-32}{3}$\end{solution}



\end{parts}
\addpoints

\question Expresa algebraicamente y simplifica los siguientes enunciados:
\begin{parts}
\part [1] La suma de dos números sabiendo que uno es 14 cm mayor que el otro

\begin{solution} $ x+x+14 = 2x+14$ \end{solution}
\part [1] El aŕea de un rectángulo cuya base mide 2 cm más que la altura

\begin{solution} $(x+2)\cdot x=x^2+2x$ \end{solution}

\part [2] El área de un triángulo sabiendo que la altura es la tercera parte de la base
\begin{solution}$\frac{x\cdot \frac{x}{3}}{2}=\frac{x^2}{6} $ \end{solution}


\end{parts}


\question[2] El padre de Ana tiene 5 años menos que su madre y la mitad de la edad de la madre es 23. ¿Qué edad tiene el padre de Ana? (Resuelve el ejercicio con una ecuación)
\begin{solution} El padre de Ana tiene 5 años menos que su madre, es decir, el padre tiene 46 - 5 = 41 años\end{solution}

\question[2] Dado un número, la suma de su mitad, su doble y su triple es 55. ¿Qué número es? (Resuelve el ejercicio con una ecuación)
\begin{solution} el número es 10\end{solution}




\addpoints

\end{questions}

\end{document}
